\documentclass[a4paper, 12pt]{article}
\usepackage{geometry}
\geometry{left=3.17cm, right=3.17cm, top=2.54cm, bottom=2.54cm}
\usepackage{graphicx}
\usepackage{amsmath}
\usepackage{enumerate}
\usepackage{amssymb,bm}
\usepackage{graphicx}
\usepackage{amsmath}
\usepackage{braket}
\usepackage{mathtools}
\usepackage{bbm}
\usepackage{booktabs}
\usepackage{tabularx}
\usepackage[colorlinks=true,linkcolor=blue,citecolor=blue,urlcolor=blue]{hyperref}
% Required for inserting images

\title{Lattice Surgery qLDPC}
\author{Tiangang Zhou}
\date{June 2026}

\begin{document}

\maketitle

\section{Introduction}

The Pauli based computation (PBC) provide a concrete way to perform the universal quantum computation. With that architectural, the magic state injection requires the measurement of $\bar{Z}_i \otimes \bar{Z}_j$, and the logical Clifford operation on the former Hilbert space that represents the processor will change the $\bar{Z}_i $ to the general operator $\bar{P}$. It's composed by general logical Pauli operator. Hence, the ultimate goal is to evaluate the joint pauli product measurement (PPM) $\bar{P}_i \otimes \bar{Z}_j$ where $P = X, Y, Z$. Besides, it's also required to evaluate the single logical qubit PPM $X_i, Y_i, Z_i$.

You may wonder why we need lattice surgery? First, we cannot destructively measurement all the single qubit since we need to keep further use this state since the magic state injection is just the intermediate step. Second, we can measure the logical operator like the stabilizer, which means at last round we construct an ancilla and apply correponding CNOT between the data qubit and the ancilla. However, this measurement is not fault tolerant with distance 1, since any error on the ancilla will cause an error. Also the logical weight can be very high, such as 14 for Bicycle Bivariate code [[98, 6, 12]]. Therefore we need the lattice surgery to address this problem. The information of the logical measurement is distributed in low-weight stabilizer. The error is regarded as normal syndrome effect that can be decoded. Besides, $d$ round of merge process can also correct the measurement error. Finally, our construction ensures the distance of the merge code won't decrease. As a result, the logical parity measurement is the extra stabilizer parity product.

\section{Single logical PPM: gadget construction}
\subsection{Construction}

The next question is how to construct the lattice surgery? First we start with the single gadget construction, to measure the logical $\bar{X}$ as an example. Since CSS codes and $3$-term chain complexes are equivalent, we build the merged code as the \emph{mapping cone} of a chain map $f:\mathcal{A}\to\mathcal{C}$ (Ide, Gowda, Nadkarni, Dauphinais, arXiv:2410.02753, Eqs.~(12)--(13)). The original code $\mathcal{C}$ is the complex $C_2 \xrightarrow{H_X^T} C_1 \xrightarrow{H_Z} C_0$ with $C_1$ the $n$ data qubits, $C_2$ the $X$-checks, $C_0$ the $Z$-checks, $H_X \in \mathbb{F}_2^{m_X \times n}$, $H_Z \in \mathbb{F}_2^{m_Z \times n}$, and $H_Z H_X^T = 0$. The ancilla is a second $3$-term complex $\mathcal{A}: A_1 \xrightarrow{\partial_1} A_0 \xrightarrow{\partial_0} A_{-1}$ with $A_1 = S_X'$, $A_0 = Q'$, $A_{-1} = S_Z'$ and $\partial_0\partial_1 = 0$. The cone of $f$ merges the original $(Q, S_X, S_Z)$ with the ancilla $(Q', S_X', S_Z')$; its check matrices are
\begin{equation}
    \tilde{H}_X = \begin{pmatrix}
        H_X & 0 \\
        f_1^T & \partial_1^T
    \end{pmatrix},
    \qquad 
    \tilde{H}_Z=\begin{pmatrix}
        H_Z & f_0 \\
        0 & \partial_0
    \end{pmatrix}.
\end{equation}
The goal is to provide the universal construction of the chain map and ancilla boundaries $f_1, \partial_1, f_0, \partial_0$.

The gadget construction is associated with the chosen logical operator. For a logical-$\bar X$ representative with support
  $x \in \mathbb{F}_2^n$ satisfying $H_Z x = 0$. We can defined several space. First, set $V_0$ is the physical qubit that logical operator $\bar{X}$ has support
  \begin{equation}
      V_0 = \mathrm{supp}(x) = (v_0,\dots,v_{|V_0|-1})
  \end{equation}
 Second, we define the set $C_0$ to be the dual stabilizer checks ($S_Z$) that have overlap with the $V_0$
 \begin{equation}
     C_0 = \{j : (H_Z)_{j,V_0} \neq 0 \} = (c_0,\dots,c_{|C_0|-1}),
  \end{equation}
Based on these two sets, we can define the projection to the sets. For any ordered subset $S=(s_0,\dots,s_{|S|-1})$, define the projection $\pi_S \in \mathbb{F}_2^{|S|\times N}$ such that
  \begin{equation}
  (\pi_S)_{i,j} = \delta_{j,s_i},\ (\pi_S^T)_{i,j} = \delta_{i, s_j} 
  \end{equation}
  We can find the projection $\pi_S M \pi_T^{T} = M[S,T]$ for arbitary matrix $M$ where $M[S,T]$ is the numpy style index.
  
  The boundary $\partial_1: A_1 \to A_0$ records how each new $X$-stabilizer acts on the ancilla qubits---the overlap of the $Z$-stabilizers with $\bar{X}$'s support---in the edge--vertex orientation of arXiv:2410.02753 Eq.~(12): $(\partial_1)_{j,i} = (\pi_{C_0} H_Z \pi_{V_0}^{T})_{j,i} = (H_Z)_{c_j, v_i}$, with $\partial_1 \in \mathbb{F}_2^{|C_0|\times|V_0|}$, so it enters $\tilde H_X$ as $\partial_1^T$. The chain map's other component attaches the new $X$-stabilizers to the original data qubits, $f_1 = \pi_{V_0}^T$ (i.e.\ $f_1^T = \pi_{V_0} \in \mathbb{F}_2^{|V_0|\times n}$). From here we find that we have $|V_0|$ new $S_X'$ ancilla stabilizers and $|C_0|$ new $Q'$ ancilla qubits.

  Since the new ancilla qubits $Q'$ are introduced, the original $Z$-checks pick up support on them through the chain-map component $f_0 = \pi_{C_0}^T \in \mathbb{F}_2^{m_Z\times  |C_0|}$. Finally $\partial_0: A_0 \to A_{-1}$ is the gauge-fixing boundary, chosen so that $\mathcal{A}$ is a chain complex, $\partial_0 \partial_1 = 0$; equivalently $\operatorname{row}(\partial_0) = \ker(\partial_1^T) \in \mathbb{F}_2^{|C_0|}$ spans the ancilla $Z$-operators commuting with the new $X$-stabilizers, i.e.\ the cycle space of $G_{\mathrm{aux}}$. (arXiv:2410.02753 writes this gauge block $G$ in Eq.~(62)/(12) and $\partial_0$ in Eq.~(68)/(13); we use $\partial_0$ throughout.) The chain-complex relation $\partial_0 \partial_1 = 0$ is exactly the CSS commutation of the two new stabilizer families. Without introducing $\partial_0$, the logical dimension is $n+|C_0|-m_X-m_Z-|V_0|$; gauge fixing removes $r= \dim \operatorname{row} \partial_0=|C_0|-|V_0|+\dim \ker(\partial_1)$ of them, leaving the merged logical dimension $k-\dim \ker(\partial_1)$. In the low-rate surgery $\dim \ker(\partial_1)=1$, since $\ker(\partial_1) \in \mathbb{F}_2^{|V_0|}$ is the single linear dependence of the stabilizers $S_X'$ on the ancilla qubits, namely $\prod_i (S_X')_i = \bar X$.

  
  Finally, the merged check matrix is
  \begin{equation}
  \tilde H_X =
  \begin{pmatrix}
  H_X & 0 \\
  \pi_{V_0} & \pi_{V_0} H_Z^{T} \pi_{C_0}^{T}
  \end{pmatrix},
  \qquad
  \tilde H_Z =
  \begin{pmatrix}
  H_Z & \pi_{C_0}^{T} \\
  0 & \ker\left(\pi_{V_0} H_Z^{T} \pi_{C_0}^{T} \right)
  \end{pmatrix},
  \end{equation}
  with the shape $\tilde H_X \in \mathbb{F}_2^{(m_X+|V_0|)\times(n+|C_0|)}$ and $
  \tilde H_Z \in \mathbb{F}_2^{(m_Z+r)\times(n+|C_0|)}$.

  We can check it satisfies the CSS commutation relation.
  Using $\pi_{V_0}\pi_{V_0}^{T}=I_{|V_0|}$ and $\pi_{C_0}\pi_{C_0}^{T}=I_{|C_0|}$,
  \begin{equation}
  \tilde H_X\,\tilde H_Z^{T} \;=\;
  \begin{pmatrix}
  H_X H_Z^{T} & 0 \\[2pt]
  \pi_{V_0}H_Z^{T}+\pi_{V_0}H_Z^{T}\pi_{C_0}^{T}\pi_{C_0} & \pi_{V_0}H_Z^{T}\pi_{C_0}^{T} \ker(\pi_{V_0}H_Z^{T}\pi_{C_0}^{T})^{T}
  \end{pmatrix}.
  \end{equation}
  The top-left vanishes by $H_X H_Z^{T}=0$. The bottom-left vanishes because
  $\pi_{C_0}^{T}\pi_{C_0} = I_{|C_0|}$ projects $H_Z^{T}$'s columns onto $C_0$ and the
  remaining rows of $\pi_{V_0} H_Z^{T}$ are zero. It means the $I_{|C_0|}$ won't cause any effect. Therefore two $\pi_{V_0} H_Z^{T}$ is canceled in the $\mathbb{F}_2$.
  The bottom-right vanishes by definition using the property of the $\ker$.

The new $X$ stabilizer is
  \begin{equation}
 (S_X')_i = \prod_{l=0}^{n-1} X_l^{(f_1^T)_{i,l}} \prod_{k=0}^{|C_0|-1} X_{\kappa_k}^{(\partial_1^T)_{i,k}}, \quad
  i=0,\dots,|V_0|-1.
  \end{equation}

The exponent on each ancilla $\kappa_k$ in $\prod_i (S_X')_i$ is
  \begin{equation}
  \sum_{i=0}^{|V_0|-1}(\partial_1^T)_{i,k}
  = \bigl(\mathbf{1}^T \pi_{V_0} H_Z^T \pi_{C_0}^T\bigr)_k
  = (H_Z x)^T \pi_{C_0}^T \big|_k
  = 0,
  \end{equation}
  where we used $\mathbf{1}^T \pi_{V_0} = x^T$ and $H_Z x = 0$, because $x$ is the logical $\bar{X}$ operator and therefore commutes with check $H_Z$. Therefore
  \begin{equation}
  \prod_{i=0}^{|V_0|-1}(S_X')_i = \prod_{l=0}^{|V_0|-1} X_{v_l} = \bar X,
  \end{equation}
  i.e., the full product of new $X$-stabilizers reproduces the measured logical.

\subsection{Boost}

So far we have provided a concrete construction of the gadget only specify a logical operator $\bar{X}$ and the original check matrix. However, to make the merged code still has the distance $O(d)$, we have to randomly add more ancilla qubits and checks to obtain a better auxiliary graph $G_{\mathrm{aux}}$. Reading $\partial_1$ as an edge--vertex incidence matrix (rows $=$ edges $Q'$, columns $=$ vertices $S_X'$) defines an auxiliary
  graph $G_{\mathrm{aux}}$ with the dictionary: vertex $\leftrightarrow$ supported qubits of logical $\bar{X}$, $(S_X')_i$,
  edge $\leftrightarrow$ ancilla $Q'$, cycle $\leftrightarrow$ new $Z$-stabilizer $(S_Z')_l$. 
  
  If the extended graph has the cheeger boundary constant $h \geq 1$ with
\begin{equation}
  h(G_{\mathrm{aux}}) = \min_{\substack{V \subseteq V_0 \\ 0 < |V| \leq |V_0|/2}} |\partial
  V|/|V| \geq 1,
  \end{equation}
  where $\partial V \subseteq E_{\mathrm{aux}}$ is the edge cut between $V$ and $V_0 \setminus V$. If initially the condition is not satisfied, we should randomly add weight-$2$ rows to $\partial_1$, or equivalently, edges to $G_{\mathrm{aux}}$. Any $Z$-logical of the merged code is then supported on $|V|$ data qubits and
  $|\partial V| \geq |V|$ ancillas, so its weight is $\Omega(d)$ and the original distance is preserved.

  \subsection{Some understanding}

  And here is the key link to surgery: the value of the measured logical operator L is not stored in the
  data qubits at one instant---it is reconstructed from the time record of the merge stabilizer outcomes
  (their parity across the rounds). So a time-like chain spanning all $\tau$ rounds flips the reported value
  of L while leaving the stored state untouched and (in the bulk) undetected. That is exactly a time-like
  logical error: an incorrect measurement of the target operator.

\section{Joint logical PPM: universal adapter}

\subsection{Construction}

  Given two CSS codes $\mathcal{C}_l$ and $\mathcal{C}_r$ with logical-$\bar X$ representatives $x_l, x_r$,
   we measure the joint Pauli $\bar X_l \otimes \bar X_r$ in a single coherent round. Apply the
  single-gadget construction on each side to obtain $(V_0^s, C_0^s, f_1^{s}, \partial_1^{s}, f_0^s, \partial_0^{s})$ with $s \in \{l, r\}$.
  
  Introduce a shared adapter $\mathcal{A} = (a_0, \dots, a_{w-1})$ of $w$ qubits carrying the parity-check
  matrix of the canonical $w$-qubit repetition code,
  \begin{equation}
  H_R \in \mathbb{F}_2^{(w-1)\times w}, \qquad (H_R)_{c, j} = \delta_{j, c} + \delta_{j, c+1}.
  \end{equation}
  Pick ordered port subsets $\mathcal{P}_s \subseteq V_0^s$ of common size $w$ on each side, identifying
  $a_i$ with the pair $(\mathcal{P}_l[i], \mathcal{P}_r[i])$.

  \subsection{Subroutines}
  We use two certain subroutine in the graph theory to modify the extended graph
  \begin{description}
  \item[\textsc{Cellulation}.] Augments $\partial_1^{s}$ with weight-$2$ rows until the port-induced
  subgraph of $G_{\mathrm{aux}}^s$ is connected and has bounded basis-cycle length.
  \item[\textsc{SkipTree}.]     By Swaroop et al.\ Thm.~7, the input is the  the cellulated port subgraph $\partial_1^{s,\mathrm{aug}}$ and the output is a $(3,2)$-sparse
  matrix $T_s$ and permutation matrix $P_{\sigma_s}$ with
   \begin{equation}
  \underbrace{T_s}_{(w-1)\times m^s_{\mathrm{aug}}}
  \Bigg(\underbrace{\partial_1^{s,\mathrm{aug}}}_{m^s_{\mathrm{aug}}\times|V_0^s|}
  \underbrace{\pi_{\mathcal{P}_s}^T}_{|V_0^s|\times w}
  \Bigg)
  \underbrace{P_{\sigma_s}}_{w\times w}
  =
  \underbrace{H_R}_{(w-1)\times w}
  \end{equation}
  computable in $O(w + m^{\mathrm{port}}_{\mathrm{aug}})$ time. The superscript $s \in \{l, r\}$ labels the side. On side $s$, write $|V_0^s|$ for
   the size of the logical-$\bar X_s$ support and $m^s_{\mathrm{aug}} = |C_0^s| + e_s$, where $|C_0^s|$ is the number of dual checks of the original code touching $V_0^s$ and $e_s$ is the
  number of weight-$2$ rows added by cellulation/boost on side $s$. Equivalently, $m^s_{\mathrm{aug}}$
  counts all ancilla qubits introduced on side $s$, both from the original single-gadget construction and
  from the boost. The port width $w$ is shared between the two sides, while $|V_0^s|, |C_0^s|, e_s$ may
  differ.

   
  \end{description}

  \subsection{Bridge stabilizers}

  With $(T_s, \sigma_s)$ for both sides and $H_R$ on $\mathcal{A}$, the merged
  check matrices, in the qubit ordering
  $(Q_l \mid Q_r \mid Q'_l \mid Q'_r \mid \mathcal{A})$ of total size
  $N = n_l + n_r + m^l_{\mathrm{aug}} + m^r_{\mathrm{aug}} + w$, take the block
  form
  \begin{equation}
  \tilde H_X^{\mathrm{joint}} =
  \begin{pmatrix}
  H_X^l & 0 & 0 & 0 & 0 \\
  0 & H_X^r & 0 & 0 & 0 \\
  f_1^{l,T} & 0 & \partial_1^{l,\mathrm{aug},T} & 0 & \pi_{\mathcal{P}_l}^T P_{\sigma_l} \\
  0 & f_1^{r,T} & 0 & \partial_1^{r,\mathrm{aug},T} & \pi_{\mathcal{P}_r}^T P_{\sigma_r}
  \end{pmatrix},
  \end{equation}
  \begin{equation}
  \tilde H_Z^{\mathrm{joint}} =
  \begin{pmatrix}
  H_Z^l & 0 & f_0^l & 0 & 0 \\
  0 & H_Z^r & 0 & f_0^r & 0 \\
  0 & 0 & \partial_0^{l,\mathrm{aug}} & 0 & 0 \\
  0 & 0 & 0 & \partial_0^{r,\mathrm{aug}} & 0 \\
  0 & 0 & T_l & T_r & H_R
  \end{pmatrix},
  \end{equation}
  where the port-label block $\pi_{\mathcal{P}_s}^T P_{\sigma_s} \in \mathbb{F}_2^{|V_0^s|\times w}$
  attaches each port vertex $v_i \in \mathcal{P}_s$ to its single adapter qubit
  $a_{\mathrm{label}_s(i)}$.

  The bridge introduces exactly one new family of stabilizer rows, which is the
  $w-1$ rows in the last block-row of $\tilde H_Z^{\mathrm{joint}}$,
%   \begin{equation}
%   (S_Z^{\mathrm{cycle}})_c
%   = \prod_{k:(T_l)_{c,k}=1} Z_{\kappa^l_k}
%     \prod_{k:(T_r)_{c,k}=1} Z_{\kappa^r_k}
%     \prod_{j:(H_R)_{c,j}=1} Z_{a_j},
%   \quad c=0,\dots,w-2.
%   \end{equation}
  It has no support on the data qubits $Q_l\cup Q_r$.
  The bridge $X$ direction is not a new row. It is encoded in the
  $\pi_{\mathcal{P}_s}^T P_{\sigma_s}$ blocks, which attach a single adapter $X$
  to each port-vertex $S_X'^s$-row.

  \paragraph{CSS commutation.}
Expand $\tilde H_Z^{\mathrm{joint}} (\tilde H_X^{\mathrm{joint}})^T$ as a
  $5 \times 4$ block matrix indexed by the $\tilde H_X$ and $\tilde H_Z$.
  Most blocks are trivially zero. But we consider the extra term $\tilde H_Z^{\mathrm{joint}} (\tilde H_X^{\mathrm{joint}})^T[4,:]$, which involves the new introduced bridge check matrix. We find the new term is to invovles the 
  \begin{equation}
  T_s \partial_1^{s,\mathrm{aug}} \;+\; H_R\, P_{\sigma_s}^T \pi_{\mathcal{P}_s}.
  \end{equation}
  The SkipTree identity is
  \begin{equation}
  T_s \partial_1^{s,\mathrm{aug}} \pi_{\mathcal{P}_s}^T = H_R\, P_{\sigma_s}^T.
  \end{equation}
  Because
  $T_s \partial_1^{s,\mathrm{aug}}$ is already zero on non-port columns of
  $V_0^s$, multiplying both sides on the right by $\pi_{\mathcal{P}_s}$ recovers
  the full matrix,
  \begin{equation}
  T_s \partial_1^{s,\mathrm{aug}} = H_R\, P_{\sigma_s}^T \pi_{\mathcal{P}_s},
  \end{equation}
  so the two terms in the block are equal and cancel in $\mathbb{F}_2$. Therefore $\tilde H_Z^{\mathrm{joint}} (\tilde H_X^{\mathrm{joint}})^T = 0$.

  \subsection{Joint logical}
  Multiplying all $w-1$ bridge $X$-stabilizers cancels at every adapter qubit, since each $a_j$ appears in exactly
  two consecutive rows of $H_R$ and every internal $\partial_1^{s,\mathrm{aug}}$ row, leaving
  \begin{equation}
  \prod_{s=l,r}\prod_{i=0}^{|V_0^s|-1} (S_{X}'^s)_i = 
   \bar X_l \otimes \bar X_r,
  \end{equation}
  so syndrome measurement returns the value of the joint Pauli.

  \subsection{Distance}
  Let $G_{\mathrm{aux}}^s|_{\mathcal{P}_s}$ be the port-induced subgraph after cellulation. The
  merged-code distance is $\Theta(\min(d_l, d_r))$ provided each side's port subgraph satisfies
  \begin{equation}
  h(G_{\mathrm{aux}}^s|_{\mathcal{P}_s}) = \min_{\substack{V \subseteq \mathcal{P}_s \\ 0 < |V| \leq w/2}} \frac{|\partial V|}{|V|}
  \;\geq\; 1,
  \end{equation}
  achieved by randomized weight-$2$ boost rows in $\partial_1^{s}$. The adapter $\mathcal{A}$ contributes a
   fixed rep-code distance $w$, so we should choose the bridge length $w \geq \min(d_l, d_r)$ to suffice the distance requirement.

\section{Single logical mixed-Pauli PPM: system merge}

To measure a single logical operator $\bar P$ that is a mixture of $\bar X$, $\bar Y$, and
$\bar Z$ on one CSS code, we introduce \emph{no} $Y$-basis ancilla. Following Ide, Gowda,
Nadkarni, and Dauphinais (arXiv:2410.02753, \S III D, Eqs.~(63)--(68)), we run the
single-gadget construction of \S 2 twice---once for the $X$-part and once, in the dual basis,
for the $Z$-part---and \emph{merge} the few stabilizer pairs that anticommute into single
non-CSS rows. The $X/Z$ mixing then lives in a \emph{check}, not in a qubit's state, so every
ancilla stays in a CSS eigenstate.

\subsection{Setup}
Write $\bar P$ through two binary representatives $x, z \in \mathbb{F}_2^n$ (the $\bar X$- and
$\bar Z$-shadows) with $H_Z x = 0$ and $H_X z = 0$. Up to an overall phase
$\bar P \propto \bar X \bar Z$, where $\bar X = \prod_v X_v^{x_v}$ and
$\bar Z = \prod_v Z_v^{z_v}$. The physical support splits into three disjoint sets,
\begin{equation}
V_X = \mathrm{supp}(x)\setminus\mathrm{supp}(z), \quad
V_Z = \mathrm{supp}(z)\setminus\mathrm{supp}(x), \quad
W   = \mathrm{supp}(x)\cap\mathrm{supp}(z),
\end{equation}
on which $\bar P$ acts as $X$, $Z$, and $Y$ respectively. The pure logical $\bar Y$ is the
special case $V_X = V_Z = \varnothing$; a single-qubit $\bar Y$ overlap is $|W| = 1$. These three
sets are not chosen directly: they are fixed by the representatives $x, z$, which are free up to
$x \to x + (\text{$X$-stabilizer})$ and $z \to z + (\text{$Z$-stabilizer})$ (the same logical
$\bar P$ up to stabilizers and sign). In particular $W = \mathrm{supp}(x)\cap\mathrm{supp}(z)$ is
exactly the $Y$-support of $\bar P$. Choose the representatives for low weight and---for fault
distance---so that $|W|\geq 2$ with both gadget graphs cyclic, which gives genuine crossing cycles
in the merged graph $\partial_1$ of \S 4.4; the degenerate $|W|=1$ tree-join leaves
$\ker \partial_1 = \ker\partial_1^{x} \oplus \ker\partial_1^{z}$ pure-CSS and is harder to keep
distance-preserving.

\subsection{Two gadgets}
Apply \S 2 to the $\bar X$-shadow $x$ to build the $X$-system: ancilla qubits $Q'^X$ indexed by
$C_0^X = \{j : (H_Z)_{j, \mathrm{supp}(x)} \neq 0\}$, with the vertex-edge incidence
\begin{equation}
\partial_1^{x} = \pi_{\mathrm{supp}(x)}\, H_Z^T\, \pi_{C_0^X}^T
\in \mathbb{F}_2^{|\mathrm{supp}(x)| \times |C_0^X|},
\end{equation}
which is exactly the single-gadget $\partial_1$ of \S 2, and the $X$-type vertex checks (one per
$v \in \mathrm{supp}(x)$)
\begin{equation}
(S_X')_v = X_v \prod_{k} X_{\kappa^X_k}^{(\partial_1^{x})_{v,k}} .
\end{equation}
Apply the $X\leftrightarrow Z$ dual of \S 2 to the $\bar Z$-shadow $z$ to build the $Z$-system:
ancilla qubits $Q'^Z$ indexed by $C_0^Z = \{j : (H_X)_{j, \mathrm{supp}(z)} \neq 0\}$, the dual
incidence
\begin{equation}
\partial_1^{z} = \pi_{\mathrm{supp}(z)}\, H_X^T\, \pi_{C_0^Z}^T
\in \mathbb{F}_2^{|\mathrm{supp}(z)| \times |C_0^Z|},
\end{equation}
and the $Z$-type vertex checks
\begin{equation}
(S_Z')_v = Z_v \prod_{k} Z_{\kappa^Z_k}^{(\partial_1^{z})_{v,k}} .
\end{equation}
The two systems use disjoint ancillas, $Q'^X \cap Q'^Z = \varnothing$. As in \S 2, the full
products reproduce the shadows, $\prod_{v\in\mathrm{supp}(x)} (S_X')_v = \bar X$ and
$\prod_{v\in\mathrm{supp}(z)} (S_Z')_v = \bar Z$.

\subsection{The merge}
For a qubit $v \in W$ both systems carry a vertex check: $(S_X')_v$ contains $X_v$ and
$(S_Z')_v$ contains $Z_v$. These two checks share \emph{exactly} the single data qubit $v$
(their ancillas are disjoint, and each vertex check touches only that one data qubit), where
$X_v$ and $Z_v$ anticommute. They therefore cannot both sit in the stabilizer group. Replace
each such pair by its Hermitian product
\begin{equation}
y_v \;:=\; i\,(S_X')_v\,(S_Z')_v
\;=\; Y_v \prod_k X_{\kappa^X_k}^{(\partial_1^{x})_{v,k}}
            \prod_k Z_{\kappa^Z_k}^{(\partial_1^{z})_{v,k}},
\qquad v \in W,
\end{equation}
a single non-CSS check acting as $Y$ on the data qubit $v$, as $X$ on its $Q'^X$ neighbours,
and as $Z$ on its $Q'^Z$ neighbours. The factor $i$ (since $Y = iXZ$) is the overall sign that
arXiv:2410.02753 instructs us to track. For $v \in V_X$ (resp.\ $V_Z$) only the $X$-system
(resp.\ $Z$-system) has a vertex check on $v$, so those checks are untouched and stay pure
$X$- (resp.\ $Z$-) type.

\subsection{Merged check matrices}
The merge identifies the $X$- and $Z$-system vertices at every $v\in W$, gluing the two gadget
graphs into one. Its incidence matrix---the merged $\partial_1$ of arXiv:2410.02753 Eq.~(66),
with rows the vertices $V_X\sqcup W\sqcup V_Z$ and columns the edges $(Q'^X \mid Q'^Z)$---is
\begin{equation}
\partial_1 =
\begin{pmatrix}
\pi_{V_X} H_Z^T \pi_{C_0^X}^T & 0 \\
\pi_{W}   H_Z^T \pi_{C_0^X}^T & \pi_{W}   H_X^T \pi_{C_0^Z}^T \\
0 & \pi_{V_Z} H_X^T \pi_{C_0^Z}^T
\end{pmatrix}
=
\begin{pmatrix}
\partial_1^{x}|_{V_X} & 0 \\
\partial_1^{x}|_{W}   & \partial_1^{z}|_{W} \\
0 & \partial_1^{z}|_{V_Z}
\end{pmatrix},
\end{equation}
where $\partial_1^{s}|_{S}$ denotes the rows of $\partial_1^{s}$ indexed by $S$. The spurious
low-weight logicals are the cycles of this glued graph, so we take a \emph{single} cycle basis
\begin{equation}
\partial_0 = \ker \partial_1
= \bigl(\, \partial_0^{(X)} \,\big|\, \partial_0^{(Z)} \,\bigr),
\end{equation}
whose $Q'^X$-columns form the $X$-system part $\partial_0^{(X)}$ (a $Z$-type support) and whose
$Q'^Z$-columns form the $Z$-system part $\partial_0^{(Z)}$ (an $X$-type support).
\emph{These two halves are the same stabilizer row}, not two separate checks: a cycle that
crosses a merged vertex $v\in W$ runs along both $Q'^X$ and $Q'^Z$ edges, so it is a single
non-CSS row carrying $Z$ on $Q'^X$ and $X$ on $Q'^Z$. Only when $|W|=1$ does $\ker \partial_1$
collapse to the two independent CSS families $\ker\partial_1^{x}$ (pure $Z$-type) and
$\ker\partial_1^{z}$ (pure $X$-type); for $|W|\geq 2$ there are $|W|-1$ crossing cycles that are
genuinely mixed and must not be split apart---they are precisely the rows whose commutation with
the $y_v$ checks is the $W$-block of $\partial_1\partial_0^T = 0$.

In the qubit ordering $(Q \mid Q'^X \mid Q'^Z)$ of total size $N = n + |C_0^X| + |C_0^Z|$, the
merged code is the non-CSS stabilizer code with symplectic check matrix
$\tilde H = (\tilde H^{(X)} \mid \tilde H^{(Z)})$---the \S 2-convention rendering of
arXiv:2410.02753 Eq.~(68):
\begin{equation}
\tilde H =
\left(\begin{array}{ccc|ccc}
H_X       & 0                   & \pi_{C_0^Z}^T    &           &                  &                     \\
\pi_{V_X} & \partial_1^{x}|_{V_X} & 0              &           &                  &                     \\
\pi_W     & \partial_1^{x}|_{W}   & 0              & \pi_W     & 0                & \partial_1^{z}|_{W} \\
          &                     &                  & H_Z       & \pi_{C_0^X}^T    & 0                   \\
          &                     &                  & \pi_{V_Z} & 0                & \partial_1^{z}|_{V_Z} \\
          & 0                   & \partial_0^{(Z)} &           & \partial_0^{(X)} & 0
\end{array}\right).
\end{equation}
The vertical bar separates the $X$- from the $Z$-Pauli columns, each ordered $(Q, Q'^X, Q'^Z)$.
Row by row: original $H_X$ and $H_Z$, extended onto the opposite system's ancillas by
$f_0^z = \pi_{C_0^Z}^T$ and $f_0^x = \pi_{C_0^X}^T$; the pure-$X$ vertex checks on $V_X$
($f_1^{(X\setminus Z)}, \partial_1^{(X\setminus Z)}$ of Eq.~(68)) and pure-$Z$ on $V_Z$
($f_1^{(Z\setminus X)}, \partial_1^{(Z\setminus X)}$); the mixed $Y$-row family on $W$, spanning
both blocks ($f_1^{(Y|x)}, \partial_1^{(Y|x)} \mid f_1^{(Y|z)}, \partial_1^{(Y|z)}$); and the
single cycle family $\partial_0 = (\partial_0^{(Z)} \mid \partial_0^{(X)})$, likewise spanning
both blocks because it is one mixed stabilizer per row.

\subsection{Commutation and the joint logical}
In the pre-merge system the only anticommuting generator pairs are
$\{(S_X')_v, (S_Z')_v\}_{v\in W}$, each sharing the single qubit $v$; every other pair commutes
(distinct data qubits, disjoint ancillas, and the \S 2 CSS relations on each side). Merging
replaces each such pair by $y_v$, which then commutes with every surviving row and with every
other $y_{v'}$, because each deleted row's unique anticommuting partner is also deleted. The
mixed cycle rows commute with every $y_v$ by construction: the symplectic product of $y_v$ with
a cycle $\partial_0$ is exactly the $v$-th entry of $\partial_1\,\partial_0^T = 0$. The result
is a (non-CSS) stabilizer code. Multiplying the surviving vertex checks and all merged rows,
\begin{equation}
\prod_{v\in V_X} (S_X')_v \;\prod_{v\in V_Z} (S_Z')_v \;\prod_{v\in W} y_v
\;=\; i^{|W|}\, \bar X\, \bar Z \;=\; \bar P,
\end{equation}
so the syndrome record returns the eigenvalue of the mixed logical $\bar P$ with no residual
operator left over.

\subsection{Initialization and measurement}
We follow the protocol of arXiv:2410.02753 (\S III C), keeping the \S 2 structure of one ancilla
layer measured for $R \geq d$ rounds. \textbf{No qubit is ever prepared in $|Y_+\rangle$.}
\begin{enumerate}
\item Prepare the $X$-system ancillas $Q'^X$ in $|0\rangle^{\otimes|C_0^X|}$ and the $Z$-system
ancillas $Q'^Z$ in $|+\rangle^{\otimes|C_0^Z|}$. The data qubits are untouched. Every qubit is
in a CSS eigenstate; the $X/Z$ mixing lives only in the checks $y_v$ and the crossing cycles.
\item Measure all stabilizers of $\tilde H$ each round for $R \geq d$ rounds. Each mixed row
($y_v$ or a crossing cycle) is read by a syndrome ancilla with \emph{mixed couplings}: $CX$ onto
its $X$-support and $CZ$ onto its $Z$-support, exactly as one measures any non-CSS check. The
$\mathrm{obs}_0$ value is the in-circuit product of the surviving vertex checks and the $y_v$
rows, equal to $\bar P$ above.
\item Detach by measuring the ancillas in their preparation-complementary basis: $Q'^X$ in $Z$
and $Q'^Z$ in $X$. The data qubits remain encoded.
\end{enumerate}
Because the $Y$ sits on a genuine data-$Y$ qubit of $\bar P$ rather than on an adapter forced
into $|Y_+\rangle$, there is no $\prod_q Y_{a_q}$ residual, and hence none of the
destructive-readout or detector null-space machinery a $|Y_+\rangle$ adapter would require.

\paragraph{The readout stabilizer product.}
Concretely, $\mathrm{obs}_0$ is the product of \emph{exactly the new vertex/merge checks}
introduced by the gadget---and of nothing else:
\begin{equation}
\mathrm{obs}_0
\;=\; \prod_{v\in V_X}(S_X')_v \;\prod_{v\in V_Z}(S_Z')_v \;\prod_{v\in W} y_v
\;=\; i^{|W|}\,\bar X\,\bar Z \;=\; \bar P,
\end{equation}
read off the in-circuit ancilla outcomes (one $M$ record per $(S_X')_v$, one $M$ per
$(S_Z')_v$, one $MX$ per $y_v$). In the merged matrix $\tilde H$ of Eq.~(68) these are the rows
below the original $H_X$/$H_Z$ blocks: the $\pi_{V_X}$ rows, the $\pi_{V_Z}$ rows, and the mixed
$\pi_W$ rows. The original $H_X$, $H_Z$ rows and the cycle rows $\partial_0$ are \emph{not}
included. For the single-$\bar Y$ Steane example ($V_X=\{2,4\}$, $V_Z=\{1,3\}$, $W=\{5\}$) these
are the five rows $\{3,4,8,9,10\}$ of $\tilde H$, whose product on the data is
$X_2X_4Z_1Z_3Y_5$. Since $\bar Y = i\bar X\bar Z = +X_2X_4Z_1Z_3Y_5$ \emph{exactly} (the $i$
cancels $X_5Z_5 = -iY_5$), the raw $\mathrm{obs}_0$ bit equals the $\bar Y$ eigenvalue:
$\mathrm{obs}_0 = 0 \Leftrightarrow \bar Y = +1$.

\paragraph{Logical-$\bar Y$ eigenstate preparation.}
The gadget itself never prepares a $|Y_+\rangle$ ancilla (above). To \emph{certify} that the
measurement returns a definite outcome we prepare the \emph{data} in a logical-$\bar Y$
eigenstate. The correct state is the codeword
\begin{equation}
|\bar Y_+\rangle = \bar S\,|\bar X_+\rangle,
\end{equation}
\emph{not} the physical product $\bigotimes_i |Y_{i,+}\rangle$. The latter is a $+1$ eigenstate
of every single-qubit $Y_i$ but is \emph{not} a codeword: the code stabilizers are random on it,
so the bare readout $[x\,|\,z]$ is non-deterministic. On the codeword every stabilizer is $+1$,
so $[x\,|\,z] = \bar Y$ is sharp. We inject the codeword fault-tolerantly: prepare
$|+\rangle^{\otimes n}$ (so $\bar X = +1$ and the $X$-syndrome is $0$), measure the
$Z$-stabilizers, and apply the feedback Pauli $X^{R s}$ with $H_Z R = I$ to cancel the random
syndrome $s$, projecting onto the syndrome-$0$ codeword $|\bar X_+\rangle$; then apply the
logical $\bar S$.

\paragraph{Transversal $S$ realizes logical $\bar S^\dagger$.}
On a doubly-even self-dual CSS code (e.g.\ Steane), the transversal phase gate $S^{\otimes n}$
implements the \emph{logical} $\bar S^\dagger$, not $\bar S$. Conjugating the logical,
\begin{equation}
S^{\otimes n}\,\bar X\,S^{\otimes n\,\dagger}
= \prod_{v\in\mathrm{supp}(x)} iX_vZ_v
= i^{\,\mathrm{wt}(\bar X)}\,\bar X\,Z|_{\mathrm{supp}(x)},
\end{equation}
and the odd-weight phase $i^{\,\mathrm{wt}(\bar X)}$ (here $i^3 = -i$), reduced modulo the
doubly-even stabilizers, gives $-\bar Y$---the conjugation rule of $\bar S^\dagger$. So the
single-qubit identity $S|+\rangle = |Y_+\rangle$ does \emph{not} lift transversally:
$|\bar Y_+\rangle = \bar S|\bar X_+\rangle$ is realized by the physical $S^{\dagger\,\otimes n}$
(and $|\bar Y_-\rangle$ by $S^{\otimes n}$), verified numerically
($\langle\bar Y\rangle = +1$ on $S^{\dagger\,\otimes n}|\bar X_+\rangle$, $-1$ on
$S^{\otimes n}|\bar X_+\rangle$). The remaining five Pauli-basis data preparations are
$|\bar 0\rangle, |\bar 1\rangle$ ($\bar Z$ eigenstates) and $|\bar +\rangle, |\bar -\rangle$
($\bar X$ eigenstates); on all four $\bar Y$ anticommutes with the prepared logical, so the
$\bar Y$ outcome is a genuine $50/50$.

\subsection{Distance}
Boost each system's graph independently to Cheeger constant $\geq 1$ as in \S 2.2, applied to
$\partial_1^{x}$ and $\partial_1^{z}$ separately. For the merged code, decompose any error in
symplectic form as $\vec e = \vec e_X + \vec e_Z$. Since
$\mathrm{wt}(\vec e) \geq \min(\mathrm{wt}\,\vec e_X, \mathrm{wt}\,\vec e_Z)$, the \S 2
single-gadget distance bound applies to $\vec e_X$ against the $X$-system and to $\vec e_Z$
against the $Z$-system independently, preserving the original distance $d$. The distance
frequently \emph{increases}, since the $X$-type measurement raises the $Z$-distance and vice
versa (arXiv:2410.02753, Eq.~(68) ff.).

\end{document}
